% Sections/8_conclusion.tex
\section{Conclusion}

This dissertation investigated how disposable vapes, as hazardous micro-WEEE, can be identified and operationalised in UAV-oriented urban imagery. The central finding is that the problem cannot be solved by training a detector alone. Vape detection depends on three connected layers: constructing a suitable dataset, evaluating the detector rigorously, and translating detections into cleanup action.

First, the project developed the Dataset Construction Engine (DCE) to address the dataset creation barrier. The DCE brought together image upload, scraping, synthetic data creation, review, SAM2-assisted annotation, VLM-assisted metadata enrichment, audit diagnostics and COCO/YOLO export. This showed that dataset construction can be treated as an engineering workflow rather than as a disconnected set of manual tasks. The DCE is especially relevant for researchers working on niche object classes where ready-made datasets do not exist.

Second, the project created and evaluated UAVape, a dedicated vape/lighter dataset for UAV-oriented litter detection. The final real dataset contained 1,204 images, 1,328 vape boxes and 712 lighter boxes. Lighters were included as an auxiliary confuser class because they are visually similar to disposable vapes and help test whether the model has learned a meaningful vape boundary. The detector benchmark compared modern YOLO families, model scales and input resolutions under a leakage-aware split. The final selected YOLO26s model, trained using the combined real, scraped and synthetic recipe, achieved a held-out vape AP50 of 0.882 and vape AP50--95 of 0.685. This demonstrates strong vape-specific detection performance under the final test protocol.

Third, the project developed the \textit{UAVape Cleanup Planner} to explore how detections could become cleanup tasks. The prototype translated detections into evidence crops, map points, confidence scores, hazard labels, task assignment and field status updates. Stakeholder walkthrough feedback supported this direction, especially for hazard labelling, evidence/location support, status reporting and overall task clarity. However, the prototype remains early workflow evidence rather than deployment evidence. It does not yet prove live UAV integration, GPS accuracy, route optimisation, municipal integration or real-shift adoption.

The project has several limitations. The real-image dataset was UAV-oriented but not captured through a full UAV deployment. The final model was evaluated on a sealed test set, but edge deployment was not tested on UAV-compatible hardware. The workflow prototype was assessed through a structured walkthrough rather than live field use. These limitations define the next steps: collect real UAV imagery, test the model on edge devices, extend evaluation across more environmental conditions, and trial the cleanup workflow during real operations.

Overall, the contribution of this dissertation is a repeatable route from hazardous litter problem to dataset, detector and workflow. Disposable vapes were the case study, but the method is wider. The same process could support future work on batteries, sharps, e-waste fragments and other small hazardous objects in parks, towpaths, pavements and pedestrian spaces. The key conclusion is that detection only matters when it becomes actionable. UAVape shows how that translation can begin.

